% Source: https://artofproblemsolving.com/community/c3h1216032_geometry
% Stable identity: aops-topic-1216032-q01
The shortest distance from $P$ to the circle is measured along the line through the centre. That line meets the circle at two points, at distances $8$ and $8+x$ from $P$.

\begin{asy}
usepackage("amsmath");
size(5cm);
defaultpen(fontsize(10pt));
pair O=(0,0); real r=161/16; pair P=(-(r+8),0);
pair T=r*dir(180-aCos(r/(r+8))*180/pi);
draw(circle(O,r),linewidth(0.9)); draw(P--T,blue+linewidth(0.9)); draw(P--(r,0),red); dot(P);dot(T);dot(O);dot((-r,0));
label("$P$",P,W); label("$T$",T,S); label("$15$",(P+T)/2,SW); label("$8$",(P+(-r,0))/2,N); label("$x$",(0,0),S);
label("\mbox{}",(min(currentpicture,user=true)+max(currentpicture,user=true))/2,black); // keeps the text after the figure black when LaTeX typesets the labels inline
\end{asy}

By the power of the point $P$ (tangent squared equals the product of the secant segments),
\[15^2=8(8+x),\qquad 225=64+8x,\qquad x=\boxed{\dfrac{161}{8}}=20.125.\]
